\subsection{Karla 湖数据准备与积水提取}
2023 年 9 月洪灾数月后，Copernicus EMSN184 仍在监测希腊色萨利 Karla 湖周边积水，分别以 12 月 18 日 Sentinel-1 和 19 日 Sentinel-2 更新洪涝轮廓并开展暴露评估\cite{emsn184}。本研究沿用观测节奏，固定经度 $22.55^\circ$--$22.85^\circ$ E、纬度 $39.43^\circ$--$39.65^\circ$ N 为研究窗；北部核查范围为 $22.63^\circ$--$22.78^\circ$ E、$39.49^\circ$--$39.60^\circ$ N。数据时相及用途见表~\ref{tab:karla-sources}。

\begin{table}[htbp]
\centering
\caption{Karla 主案例的数据、用途与时相。官方轮廓不参与分割或阈值选择。}
\label{tab:karla-sources}
\small
\begin{tabularx}{\linewidth}{@{}>{\raggedright\arraybackslash}p{32mm}>{\raggedright\arraybackslash}p{34mm}>{\raggedright\arraybackslash}X@{}}
\toprule
来源 & 时相 & 用途\\
\midrule
Sentinel-1 RTC & 2023-09-01；2023-12-18 & 同相对轨道 7 的 VV/VH；提供灾前基线和首期雷达候选。\\
Sentinel-2 L2A & 2023-08-31；2023-12-19 & 绿光、短波红外和场景分类；排除不可判读像元，独立提取光学候选。\\
EMSN184 & 2023-12-18；2023-12-19 & 官方作用域及两期轮廓，仅用于分割后的日期配对对照。\\
ESA WorldCover & 2021 & 用类别 40 表示耕地，供空间交叠统计。\\
OpenStreetMap & 2026-11-25 获取 & 主要道路及连接道快照，供现状定位。\\
\bottomrule
\end{tabularx}
\end{table}

下载工具先向 Microsoft Planetary Computer 的 STAC（时空资产目录）提交研究窗、传感器和观测日期，按轨道和所需波段筛选场景，再用 GDAL 从云端读取研究窗内的像元\cite{planetary}。目录记录提供的是观测足迹和波段地址；实际可判读覆盖率在读取像元后计算。两类影像统一到以米为单位的 UTM 34N 投影（EPSG:32634）和 20 m 网格，保证同一行列代表同一位置，每格对应 400~$\mathrm{m}^2$。

雷达工具读取 VV 和 VH 两种极化的后向散射强度，并换算为 dB。光学工具读取绿光 B03、短波红外 B11 和场景分类图层（SCL）；反射率波段采用双线性重采样，离散的场景类别保留最近邻值。同一像元所需的波段取自同一场景，拼接时优先保留可判读的像元。下载清单记录场景编号、观测日期、有效覆盖率及输出文件的 SHA-256 校验值。

雷达变化法对 VV/VH 做 $3\times3$ 中值处理，要求灾后 VV $<-16$ dB、VH $<-23$ dB，且 VV 相对灾前下降超过 3 dB，排除灾前已满足双极化暗阈值的区域。光学方法以绿光 $G$ 与短波红外 $S$ 构成改进型归一化差异水体指数 $\mathrm{MNDWI}=(G-S)/(G+S)$\cite{mndwi}；纳入 SCL 类别 4、5、6 且反射率和超过 0.015 的像元。MNDWI 超过 0.12 的灾后区域扣除灾前检出的水体。两种方法均剔除少于 9 个像元的八连通斑块。

分析程序先检查灾前与灾后像元是否对齐，再计算两期均可判读的像元占比；低于 65\% 时不发布该次分割。输出包括标记候选积水、灾前水体、有效但未检出和不可判读的 GeoTIFF 掩膜，以及由候选斑块生成的 GeoPackage 矢量图层。发布前核对多边形面积与候选像元数乘以单格面积所得结果。

比较时，仅在共同可判读区内判定两期均检出、仅雷达检出、仅光学检出和均未检出；区外标为不可比。工具再与同日官方轮廓对照，按官方作用域内的有效像元计算交并比。暴露核查按证据类别计数落在耕地内的像元，用像元数乘以单格面积得到耕地重叠面积；道路则先裁到统计范围，去除重合线段，再与各类别范围求交，把交叠线段按类别累加。两种结果都保留原始证据类别，便于区分一致检出、单一传感器检出和不可比范围。
